% Kapitel 7 - Offline Domain-Gap Evaluation
label{sec:domaingap}
\section{Offline Domain-Gap Evaluation}

\paragraph{How to interpret the results.}
The results in this chapter quantify structural discrepancies rather than visual realism or driving performance.
Observed differences should be interpreted as indicative of representational limitations in automated map generation pipelines, not as direct measures of simulator realism or model safety.

\subsection{Outcome of manual vs.\ auto tile correspondence}
\label{sec:correspondence_outcome}
For the authoritative run\_11 artifact, spatial correspondence construction produced zero valid matches (\texttt{correspondence.csv} contained only the header), and results were therefore limited to whole-map structural metrics.

A subsequent analysis (run\_14) using the full-pipeline-processed auto XODR in the canonical CRS produced \textbf{39 tile pairs} with spatial IoU $\geq 0.5$; all 39 pairs had IoU $= 1.0$ and 27 passed the IoU gate (mean IoU over all pairs $= 0.692$). This confirms that the zero-correspondence outcome in run\_11 was a CRS artefact of using an unprocessed auto XODR: once the auto map is projected into the same canonical transverse-Mercator frame as the manual tiles, per-tile correspondence is fully resolved.

The per-tile gap metrics from run\_15 provide the first complete localized domain-gap evidence for this thesis: the run yields \textbf{27 tile pairs} all with IoU $= 1.0$ (of 39 total, 27 passing the IoU $\geq 0.5$ gate, mean IoU $= 0.692$). Over these 27 tiles, the per-tile geometry RMSE has a mean of $46.9$\,m and a median of $36.0$\,m; the per-tile Hausdorff has a mean of $503$\,m and a maximum of $1{,}280$\,m (worst tile: \texttt{tile\_8\_2}). The substantially lower per-tile mean Hausdorff ($503$\,m) compared with the whole-map figure ($1{,}526$\,m) confirms that the whole-map metric was dominated by the large coverage mismatch between the auto and manual networks rather than by local geometry error within overlapping tiles. The earlier run\_14 single-tile result (tile\_9\_1: RMSE $= 98.9$\,m, Hausdorff $= 261$\,m) is superseded by these aggregate statistics.

\label{sec:domaingap}

This chapter defines the offline protocol used to quantify differences between automatically generated maps and the manual OpenDRIVE reference. Offline evaluation is designed to remain meaningful even if CARLA ingestion fails for some tiles.

\subsection{Reference--candidate contract}
Domain-gap evaluation compares:
\begin{itemize}
\item \textbf{Reference maps:} tiles extracted from the manually authored OpenDRIVE map.
\end{itemize}
The reference map is treated as immutable; only deterministic preprocessing (tiling, metadata extraction) is permitted.

\paragraph{Pre-processing pipeline for candidate maps.}
\label{par:sumo_preprocessing}
Before domain-gap analysis, the raw osm2xodr output passes through two deterministic preprocessing stages that are part of the governed pipeline and apply equally to all runs reported in this thesis.
\textit{Stage 1 (sanitization)} removes malformed XML, normalises road identifiers, and enforces OpenDRIVE schema compliance without altering road geometry.
\textit{Stage 2 (SUMO netconvert topology repair)} passes the sanitized XODR through SUMO netconvert 1.24.0 (flags: \texttt{--geometry.remove true --ignore-errors}) to reconstruct junction connectivity and remove collinear intermediate waypoints.
The \texttt{--geometry.remove} flag merges consecutive geometry segments whose shared waypoint lies within a configurable angular tolerance, reducing vertex count without altering the road centreline at a macroscopic scale.
All geometry and curvature metrics reported in this chapter are therefore computed on the candidate XODR that has passed through Stage~2 (SUMO topology repair) and the subsequent enrichment stages (DEM elevation injection, lane attribute normalisation, building and signal enrichment — Stages~3--8). The Stage-2 \texttt{--geometry.remove} simplification is thus one of several deterministic transformations applied before metric computation; it is the most relevant to curvature sampling because it alters waypoint density.
\textit{Stage 4 (OSM metadata enrichment)} additionally derives lane widths from the RASt~06 road-type catalogue, writes \texttt{<speed>} elements from the OSM \texttt{maxspeed} tag, and stamps turn-lane hints (\texttt{userData/turnMarking}) from the OSM \texttt{turn:lanes} tag where present. German regulatory signs (\textit{e.g.}\ StVO~206 stop, 274-xx speed-limit) are injected as \texttt{<object>} elements for roads with an explicit \texttt{traffic\_sign} OSM tag. These metadata enrichments operate on road attributes and do not alter road centreline geometry; they therefore have no effect on the primary structural metrics (RMSE, Hausdorff, curvature divergence) reported below, but improve semantic fidelity for downstream CARLA simulation.
Post-repair quality gates check junction integrity and lane-link target existence; if SUMO worsens either count, the pipeline raises an error and falls back to the sanitized input.
Logged SUMO diagnostics for the run\_11 candidate confirm: 5 no-lane edges, 97 sharp-turn warnings (radius $< 6$\,m), 84 geometrically incompatible lane connections, and 43 stop-line misalignments; these are carried forward as known structural deficiencies rather than suppressed.

\subsection{Geometric alignment}
Because global offsets may exist, alignment is performed prior to metric computation. A similarity transform in 2D is estimated:
\begin{equation}
\mathbf{x}' = s \mathbf{R}\mathbf{x} + \mathbf{t}, \qquad
\mathbf{R} \in SO(2),
\label{eq:sim_transform}
\end{equation}
where $s$ is scale, $\mathbf{R}$ is rotation, and $\mathbf{t}$ is translation. If needed, ICP refinement minimizes point-to-curve distances \cite{besl1992icp}.

\subsection{Structural metrics}
After alignment, the following complementary metrics are computed.

\paragraph{Pointwise RMSE.}
Given $N$ sampled points $\{\mathbf{p}_i\}$ from the candidate and their closest points $\{\mathbf{q}_i\}$ on the reference,
\begin{equation}
\mathrm{RMSE} = \sqrt{\frac{1}{N}\sum_{i=1}^{N}\lVert \mathbf{p}_i - \mathbf{q}_i \rVert_2^2}.
\label{eq:rmse}
\end{equation}

\paragraph{Hausdorff distance.}
Worst-case deviation is captured by the (directed) Hausdorff distance:
\begin{equation}
d_H(P,Q)=\max_{p\in P}\min_{q\in Q}\lVert p-q\rVert_2 .
\label{eq:hausdorff}
\end{equation}

\begin{figure}[ht]
\centering
\includegraphics[width=0.9\linewidth]{../images/fig_rmse_hausdorff_illustration.png}
\caption{Illustration of average vs.\ worst-case deviation: RMSE captures average drift, while Hausdorff emphasizes localized outliers.}
\label{fig:rmse_hausdorff}
\end{figure}

\paragraph{Curvature distribution divergence.}
Curvature is sampled along lanes and binned into histograms $p(k)$ and $q(k)$. The implemented divergence metric in the governed pipeline is the Kullback--Leibler divergence:
\begin{equation}
\mathrm{KL}(p\Vert q)=\sum_i p_i \log \frac{p_i}{q_i},
\label{eq:kl_curvature}
\end{equation}
with numerical stabilization to avoid empty-bin singularities. In the final thesis interpretation, this KL value is reported together with the curvature standard deviations because the distributions are strongly dominated by near-straight segments.

\subsection{Observed structural variability under fixed inputs}
A central question of this thesis is whether repeated map generation from fixed OSM inputs introduces measurable structural variability without explicit appearance randomization.
Two experiment classes are treated separately.

\paragraph{Determinism baseline (same OSM, same settings).}
Starting from a fixed OSM snapshot and a fixed pipeline configuration (including seeds where applicable), the OSM$\rightarrow$OpenDRIVE$\rightarrow$CARLA conversion is repeated $N$ times.
Outputs are compared using (i) file hashes of the produced \texttt{.xodr} files, and (ii) structural statistics (road count, lane count, junction count, total road length).
This experiment addresses determinism only; it does not itself constitute domain randomization.

\paragraph{Ensemble generation (many maps, fixed geography).}
Variability analysis is treated as a distinct experiment class in which multiple maps are generated for the same Ingolstadt bounding box under controlled parameter changes or repeated end-to-end generation.
This is conceptually related to domain randomization \cite{tobin2017domainrandomization}, but the thesis does not claim intentional appearance randomization in the sense of texture, lighting, or weather perturbation. Earlier repository notes sometimes used the informal phrase ``natural domain randomization'' for this phenomenon; here it is described more precisely as observed structural variability under fixed inputs while keeping the geographic extent fixed.

\subsection{Perceptual and task-level evaluation (RQ3--RQ5)}
Direct paired capture between manual and generated Ingolstadt maps is a prerequisite for perceptual-gap quantification.
The governed pipeline sensor rig (6 cameras, 2 LiDARs) was validated on the built-in \texttt{Town10HD} map, successfully recording \textbf{20 RGB frames} with all 14 listener callbacks registered (\texttt{ok=true}); T-PERC-LISTENER-ASSERT-001 is closed. Capture on the Ingolstadt cooked maps (\texttt{Grid0821} and \texttt{Grid0828}) is blocked by a confirmed CARLA~0.9.16 resource boundary (T-PERC-GRID0821-SPAWN-CRASH-001, closed/deferred): \texttt{client.load\_world(`Grid0821')} itself crashes the CARLA process on this hardware, indicating that the map's cooked asset size exceeds available GPU VRAM or system RAM. This is not a defect in the governed pipeline code. RQ3 evidence is therefore bounded at rig-code correctness on \texttt{Town10HD}; a full paired Ingolstadt comparison is classified as future work. Earlier development-run references to 241 frames are not used as primary thesis evidence in this governed tree.
The governed GNN artifacts are also bounded: the earlier positive-only cosine-loss interpretation was superseded by a contrastive run with \texttt{verdict = NOT\_COLLAPSED}, but no downstream transfer experiment was executed and no transfer-feasibility claim is made.

\subsection{Summary}
The offline domain-gap framework supports robust whole-map structural evaluation even when simulator execution is unstable or per-tile correspondence is invalid.

\section{Quantitative Results}
\label{sec:structural_results}

\subsection{RQ1: Determinism and Natural Randomization}
The final governed determinism authority is the \texttt{scenario\_b\_audit/evidence/determinism} artifact family. Its 5-run audit, executed with identical inputs, yielded a verdict of \textbf{NONDETERMINISTIC}.
While topological metrics (road counts, junction counts) remained constant ($CV=0.0$), binary file hashes differed across all runs.
This is reported as a bounded observational variability finding; no robustness or performance-improvement claim is made from this result alone.

\begin{table}[h]
\centering
\begin{tabular}{lrrrrr}
\hline
Run Index & 0 & 1 & 2 & 3 & 4 \\
\hline
Status & PASS & PASS & PASS & PASS & PASS \\
Tile Count & 75 & 75 & 75 & 75 & 75 \\
Road Count & 5,712 & 5,712 & 5,712 & 5,712 & 5,712 \\
Junction Count & 680 & 680 & 680 & 680 & 680 \\
XODR Hash & \texttt{ac6d...} & \texttt{3542...} & \texttt{e813...} & \texttt{e4c0...} & \texttt{b24d...} \\
\hline
\end{tabular}
\caption{Final governed determinism audit results (\texttt{scenario\_b\_audit}, 5 runs). Structural counts are invariant ($CV=0.0$), while cryptographic hashes diverge across all runs. The determinism audit was executed under the RQ1 pipeline configuration (5,712 roads per run; $CV = 0.0$). The structural overview in Table~\ref{tab:env_comparison} reports 5,837 roads from the RQ2 governed run, reflecting a different pipeline configuration used for the structural gap analysis. Both counts are internally consistent within their respective experiment classes.}
\label{tab:determinism_results}
\end{table}

\subsection{RQ2: Structural Domain Gap}
Comparison between the imported run\_11 auto artifact and the manually authored Grid0828 reference yields the following primary planar structural results.

\begin{table}[h]
\centering
\begin{tabular}{lrrp{0.29\linewidth}}
\hline
Metric & Auto / value & Manual / reference & Interpretation \\
\hline
Total road length & 258.6 km & 52.9 km & Coverage ratio $= 4.88\times$ \\
Junction count & 779 & 119 & Junction ratio $= 6.55\times$ \\
Whole-network RMSE & 54.84 m (frozen) & --- & Primary run\_11 planar geometry metric \\
Matched RMSE & 24.47 m (frozen) & --- & Defined on the 73.0\% matched subset within a 50\,m criterion \\
Hausdorff distance & 1662.99 m (frozen) & --- & Whole-network worst-case deviation \\
Corrected curvature KL divergence & 0.00031 (frozen) & --- & Corrected support measurement on the authoritative run\_11 pair; see note below \\
Curvature std & 0.865\,m$^{-1}$ & 0.055\,m$^{-1}$ & Spread of absolute curvature values across all road segments \\
Predecessor link rate & 0.0\% & 93.4\% & 93.4\,pp gap; auto map (contract-run) declares no explicit road predecessors \\
Successor link rate & 0.0\% & 96.7\% & 96.7\,pp gap; auto map relies on junction-based routing only \\
\hline
\multicolumn{4}{l}{\textit{Per-tile metrics (run\_15, 27 tile pairs with IoU $\geq$ 0.5, mean IoU $= 0.692$; aggregate over all matched tiles)}} \\
Per-tile RMSE (mean) & 46.9 m & --- & Mean over 27 matched tiles; run\_15 (median: 36.0 m) \\
Per-tile RMSE (max) & 118.3 m & --- & Worst-case tile (\texttt{tile\_8\_2}); run\_15 \\
Per-tile Hausdorff (mean) & 503 m & --- & Mean over 27 matched tiles; 3.0$\times$ lower than whole-map Hausdorff \\
Per-tile Hausdorff (max) & 1{,}280 m & --- & Worst-case tile; run\_15 \\
\hline
\end{tabular}
\caption{Primary structural domain-gap metrics. Geometry and curvature results are from the governed run\_11 raw basis and are 2D planView metrics only; the road-length ratio uses the audited run\_11 coverage context; junction counts use the governed connectivity companion artifact family; curvature statistics use the corrected support measurement on the same authoritative manual-vs-contract-run pair, recorded in \texttt{CURVATURE\_CLAIM\_BASIS.json}. Connectivity link rates are from the governed run\_12 companion artifact (contract-run XODR, 5{,}837 roads). Per-tile metrics are aggregate statistics over the 27 IoU-matched tile pairs from run\_15 (full-pipeline elevated+enriched XODR, 39 pairs generated, 27 passing IoU gate, mean IoU $= 0.692$). The lower per-tile mean Hausdorff ($503$\,m) compared to the whole-map value confirms domination by coverage mismatch.}
\label{tab:structural_gap}
\end{table}

\paragraph{Curvature interpretation.}
\label{par:curvature_interp}
The corrected curvature support measurement on the authoritative run\_11 pair, recorded in \texttt{CURVATURE\_CLAIM\_BASIS.json}, replaces the earlier value of $\mathrm{KL}=2493.53$, which was an artifact of the measurement code ignoring \texttt{paramPoly3} geometry elements.
Run_15 was executed with CRS alignment disabled for the whole-map stage, yielding a run_15 whole-map RMSE of $66.4$,m (not reported as a primary metric); the authoritative whole-map authority remains run_11 at $54.84$,m (CRS alignment applied). Per-tile metrics from run_15 are independent of this difference because per-tile alignment uses tile-local correspondence.
The auto XODR is composed of 71\% \texttt{paramPoly3} segments; the earlier extractor read only \texttt{arc} elements and therefore assigned $\sigma_{\mathrm{auto}}=0$, producing a degenerate KL divergence.
After applying the corrected sampler (Frenet--Serret formula: $\kappa = |U'V'' - V'U''| / (U'^2+V'^2)^{3/2}$), both distributions become non-degenerate.
The near-zero corrected KL ($0.00031$) reflects that both histograms are dominated by straight segments ($\kappa \approx 0$) when straight-segment samples are included.
The scientifically informative quantity is the standard deviation gap: $\sigma_{\mathrm{auto}} = 0.865\,\mathrm{m}^{-1}$ versus $\sigma_{\mathrm{manual}} = 0.055\,\mathrm{m}^{-1}$.
This indicates that the procedurally generated map contains a substantially wider spread of curvature values, including high-curvature outlier segments, compared with the smoothly hand-authored reference.

\begin{figure}[h]
\centering
\includegraphics[width=0.96\linewidth]{../images/fig_curvature_distribution_comparison.png}
\caption{Corrected paramPoly3-aware curvature distributions for the authoritative Ingolstadt pair (manual Grid0828 reference versus contract-run auto XODR), with straight segments included exactly as in the authoritative run\_11 measurement pair (see \texttt{CURVATURE\_CLAIM\_BASIS.json}). Left: both distributions are strongly concentrated near $\kappa \approx 0$, which explains why the corrected KL divergence remains near zero. Right: on a log-density scale, the auto map shows a much heavier high-curvature tail and rare outliers, consistent with $\sigma_{\mathrm{auto}} = 0.865\,\mathrm{m}^{-1}$ versus $\sigma_{\mathrm{manual}} = 0.055\,\mathrm{m}^{-1}$. The plotted x-ranges are cropped for readability; rarer auto-map outliers extend beyond the visible range.}
\label{fig:curvature_distribution_comparison}
\end{figure}

Figure~\ref{fig:curvature_distribution_comparison} makes this interpretation visually explicit. The near-zero KL should therefore not be read as ``curvature similarity'' in a broad sense; for thesis interpretation, the decisive signal is the much wider auto-map spread and its heavier tail of high-curvature segments.

\begin{figure}[h]
\centering
\includegraphics[width=0.86\linewidth]{../images/fig_overlay_manual_vs_auto.png}
\caption{Visual overlay of the auto-generated and manually authored Ingolstadt map extents after reprojection and rigid alignment. The figure illustrates the large coverage mismatch that contextualizes the whole-network RMSE and Hausdorff metrics; it should not be read as proof of uniformly good local alignment everywhere in the city.}
\label{fig:overlay_manual_vs_auto}
\end{figure}

\paragraph{Regression verification (run\_17).}
A post-fix regression run (run\_17, 2026-04-05) using the same canonical contract XODR as run\_11, but with the updated alignment implementation (rigid-scale-lock correction applied to \texttt{geo\_alignment.py}), yielded a whole-network RMSE of $53.93$\,m ($\Delta = -0.92$\,m, $-1.7\%$ relative to run\_11) with \texttt{parity\_check.ok=true} and 85 tests passing. This confirms that the primary run\_11 measurement is stable under minor alignment improvements. The run\_11 value of $54.84$\,m remains the frozen authoritative thesis value; run\_17 is the regression verification only.

\paragraph{Geometry and ICP interpretation.}
\label{sec:crs_normalization}
The imported run\_11 alignment artifacts show that the auto map was first reprojected from its local transverse Mercator frame into EPSG:32632 before rigid alignment. This reduced the centroid separation from 5,444,300.68 m to 760.95 m and increased plan-view bounding-box IoU to 0.637, which is the main practical improvement that made a meaningful comparison possible.
After this CRS-normalized centroid-based initialization, an ICP refinement step was executed for 20 iterations. However, the logged diagnostic fit metric is explicitly a \emph{planView geometry start-point correspondence RMSE}, and that reported metric worsened monotonically from 2060.96 to 2066.98. Because the logged diagnostic is not proven equivalent to the ICP optimization objective, the thesis does not claim successful ICP convergence. Run\_11 is therefore reported only as the best obtained rigid alignment under the evaluated setup, with the effective usable improvement coming from CRS reprojection plus centroid initialization.

\begin{table}[h]
\centering
\begin{tabular}{lp{0.20\linewidth}p{0.42\linewidth}}
\hline
Alignment stage & Value & Interpretation \\
\hline
Centroid delta before CRS reprojection & 5,444,300.68 m & Coordinate-frame mismatch dominates the unnormalized comparison \\
Centroid delta after CRS reprojection & 760.95 m & Practical whole-map placement improvement before iterative refinement \\
Run\_11 whole-network RMSE & 54.84 m & Final primary structural gap metric at whole-network scale \\
\hline
\end{tabular}
\caption{Alignment interpretation using the governed run\_11 raw artifacts. Historical pre-alignment RMSE values such as 736.65 m are retained only as earlier workflow context and are not used as the final thesis result.}
\label{tab:two_component_gap}
\end{table}

\begin{figure}[h]
\centering
\includegraphics[width=0.95\linewidth]{../images/fig_crs_normalization.png}
\caption{CRS normalization workflow used to place auto-generated and manual OpenDRIVE maps into a common metric frame before structural comparison.}
\label{fig:rq2_crs_normalization}
\end{figure}

\begin{figure}[h]
\centering
\includegraphics[width=0.95\linewidth]{../images/fig_tile_correspondence.png}
\caption{Tile correspondence diagnostics. Per-tile pairing is evaluated with a shared grid and IoU gating, while whole-map metrics remain the authoritative RQ2 evidence when tile confidence is low.}
\label{fig:rq2_tile_correspondence}
\end{figure}

\paragraph{DEM and elevation status.}
\label{par:dem_status}
The original run\_11 DEM coverage gate passed at 92.7\%, but the associated auto XODR retained zero-valued elevation coefficients throughout; all run\_11 structural metrics are therefore planar-only.
A subsequent pipeline correction introduced an expanded DEM (\texttt{dem\_ing\_expanded.tif}) covering 100\% of the road network, together with a three-tier elevation fallback system: (i) direct DEM sampling with a five-point neighbourhood probe; (ii) inverse-distance-weighted extrapolation via a \texttt{scipy} k-d tree using the five nearest valid road anchors; and (iii) road-graph BFS propagation through predecessor/successor links.
The governed run\_13 rerun confirmed \texttt{dem\_nodata\_ratio = 0.0} and \texttt{parity\_check.ok = true}; the aligned auto XODR in that run carries non-flat elevation profiles in the Ingolstadt range of 362--392\,m, with no z\,=\,0 roads remaining.
An elevation-gap comparison between the auto and manual maps was not available from run\_13 because the primary input XODR fed to the domain-gap stage was the baseline flat file.
The governed run\_15 full-pipeline rerun subsequently provided this evidence.
Using a polyline minimum-distance matcher with a 200\,m threshold, the elevation gap component matched \textbf{743 auto/manual road pairs} (of 5{,}712 auto and 972 manual roads), yielding an elevation profile RMSE of \textbf{4.40\,m} (per-road p50 $= 3.05$\,m, p95 $= 8.90$\,m; mean match distance $= 45.5$\,m).
Of the 5{,}712 auto roads, 20.9\% still carry flat (\texttt{z = 0}) elevation profiles in the run\_15 input XODR; all manual roads carry non-flat profiles.
This elevation gap result is marked \textit{supplementary} in run\_15 because the primary structural-gap authority (run\_11) remains planar; the 4.40\,m figure provides a directional bound on the elevation component of the domain gap rather than a paired planar-equivalent comparison.

\begin{table}[h]
\centering
\begin{tabular}{lrr}
\hline
Object class & Auto artifact & Manual reference \\
\hline
Barriers & 0 & 1,538 \\
Buildings & 0 & 992 \\
Poles & 0 & 3,427 \\
Trees & 0 & 4,232 \\
\hline
\end{tabular}
The contract-run XODR contains 1{,}363 	exttt{<object>} elements overall (road signs and speed-limit entries injected by the topology stage); these are not counted in the four categories above, which target 3D environmental objects only.}
\label{tab:object_gap}
\end{table}

\paragraph{Building injection fix (supplementary evidence).}
\label{par:building_injection}
The zero-building count in the contract-run artifact reflects the state of the pipeline at the time of the authoritative run\_11 measurement.
A subsequent fix addressed two issues: the \texttt{BuildingExtruder} was inserting \texttt{<object>} elements at the OpenDRIVE root level rather than inside individual \texttt{<road><objects>} elements (incorrect schema placement), and the \texttt{buildings.geojson} source file contained zero features.
Stage-4 supplementary evidence confirms that after these fixes, 5{,}683 building objects and 3{,}854 lamp-post objects are correctly injected into the intermediate XODR.
The governed run\_13 artifact carries 3{,}854 lamp-post objects in the auto XODR (confirming lamp injection is end-to-end functional) but does not include buildings, because the supplementary elevated run from which run\_13 derives its pre-anchored input was executed with building injection disabled.
The governed run\_15 full-pipeline rerun with building injection enabled confirms \textbf{5{,}683 building objects} and \textbf{3{,}854 lamp-post objects} correctly injected into the enriched auto XODR, matching the stage-4 supplementary evidence count.
The semantic-object gap in Table~\ref{tab:object_gap} therefore remains accurate for the contract-run (run\_11) artifact; the building gap is reversed in the fully enriched run\_15 context (auto: 5{,}683 buildings $>$ manual: 992 buildings), reflecting the higher coverage area of the auto map.

\paragraph{Companion structural descriptors.}
The governed companion artifact set reports: normalized intersection gap $= 0.1372$ (run\_15), normalized semantic gap $= 0.1662$, normalized road-class gap $= 0.0597$, and the connectivity link-declaration gap from run\_12 (contract-run XODR): predecessor gap $= -93.4$\,pp, successor gap $= -96.7$\,pp.
The run\_15 full-pipeline rerun with the link-patched XODR shows higher declaration rates (1.0\% predecessors, 27.8\% successors) due to the \texttt{lane\_successor\_fixed} and \texttt{linkpatched} processing stages applied to the input; this demonstrates a partial reduction in the connectivity gap when post-processing fixes are applied, but does not change the contract-run measurement.
The run\_15 composite domain-gap score is $0.637$ on the $[0,1]$ scale (components: geometry $\times 0.35$ + curvature $\times 0.20$; higher = greater gap).
These descriptors supplement the run\_11 geometry metrics but do not replace the primary whole-network RMSE, matched RMSE, and Hausdorff results reported above.

\begin{figure}[h]
\centering
\includegraphics[width=0.96\linewidth]{../../../thesis_results/structural_gap_v1/run_11/fig_road_type_distribution.png}
\caption{Road-type distribution comparison between the auto-generated and manually authored Ingolstadt maps (run\_15). The auto map shows a substantially higher proportion of residential and unclassified road segments relative to the manually authored reference, contributing to the normalized road-class gap of $0.0597$.}
\label{fig:road_type_distribution}
\end{figure}

\begin{figure}[h]
\centering
\includegraphics[width=0.96\linewidth]{../../../thesis_results/structural_gap_v1/run_11/fig_intersection_density.png}
\caption{Intersection density comparison between the auto-generated ($779$ junctions) and manually authored ($119$ junctions) Ingolstadt maps (run\_11). The auto map's $6.55\times$ higher junction count reflects the over-segmented OSM topology that results in numerous implicit intersections in procedural generation, contributing to the normalized intersection gap of $0.1372$ (run\_15).}
\label{fig:intersection_density}
\end{figure}

\begin{figure}[h]
\centering
\includegraphics[width=0.96\linewidth]{../../../thesis_results/structural_gap_v1/run_11/fig_semantic_depletion_bars.png}
\caption{Semantic-object gap across four object categories (barriers, buildings, poles, trees) for the contract-run auto XODR versus the manual Grid0828 reference. The auto contract-run artifact contains zero objects in all categories prior to the building-injection fix; the figure contextualizes the normalized semantic gap of $0.1662$ (run\_15).}
\label{fig:semantic_depletion}
\end{figure}

\subsection{Perception Capability and Readiness}
Perception readiness was evaluated by executing the full sensor rig (6 cameras, 2 LiDARs) in synchronous mode.
A batch of seven code-level defects was identified and resolved (commits \texttt{4b0ab1f3}, \texttt{d765fb7f}, 2026-04-05): synchronous mode was reordered to precede sensor listener registration; frame accounting was corrected to use the CARLA \texttt{WorldSnapshot.frame} identifier rather than a local counter; disk I/O was moved to a background executor to prevent callback blocking; ego spawning was hardened from \texttt{try\_spawn\_actor()} to \texttt{spawn\_actor()} with explicit exception handling; and manifest writing was deferred until all pending callbacks drain.

\paragraph{Town10HD full-rig capture (resolved).}
Following resolution of the \texttt{libcarla} streaming assertion (T-PERC-LISTENER-ASSERT-001) and correction of the frame-count minimum gate to 8 frames, the complete 8-sensor rig was executed on the built-in \texttt{Town10HD} map. The run recorded \textbf{20 RGB frames} with all 14 listener callbacks registered (\texttt{ok=true}, 2026-04-14). This confirms that the synchronous-mode ordering, background I/O executor, exception-path listener cleanup, and sensor-destroy path are all correct. T-PERC-LISTENER-ASSERT-001 is closed as resolved.

\paragraph{Grid0821 — confirmed CARLA resource boundary (T-PERC-GRID0821-SPAWN-CRASH-001, closed).}
Capture on the cooked Ingolstadt map \texttt{Grid0821} was investigated through three successive fix iterations.
The batch \texttt{apply\_batch\_sync()} call was replaced by sequential \texttt{world.try\_spawn\_actor()} calls (commit~\texttt{9feaaac4}), which confirmed that the first sensor actor spawns and fires its first callback; the CARLA process then hangs (\texttt{record\_route\_returncode\,=\,124}) under the full 8-sensor rig load.
A further iteration (commit~\texttt{c786bd21}) added \texttt{world.tick()} synchronisation between spawns and introduced a 2-sensor reduced rig (env\,\texttt{UP\_THESIS\_REDUCED\_RIG=1}) as a diagnostic tool; however, during the probe, a \texttt{client.load\_world(`Grid0821')} call itself caused the CARLA~0.9.16 process to terminate, and the RPC port became permanently unresponsive.
This establishes that the failure occurs at the map-loading level, not only during sensor allocation: Grid0821 exhausts available GPU VRAM or system RAM when CARLA attempts to load its cooked assets on this hardware configuration.
\textbf{T-PERC-GRID0821-SPAWN-CRASH-001 is therefore classified as a CARLA~0.9.16 runtime boundary} on the available hardware, not a defect in the governed pipeline code.
The sensor rig is confirmed correct on \texttt{Town10HD} (20 frames, \texttt{ok=true}); a paired Ingolstadt capture on Grid0821 or Grid0828 is deferred to future work requiring either higher-resource hardware or a newer CARLA version.

\paragraph{RQ3 perception status.}
The full 8-sensor rig is confirmed correct on \texttt{Town10HD} (20 frames, ok=true). Paired Ingolstadt perceptual-gap quantification remains bounded: perceptual gap measurement requires synchronized manual/simulation perceptual pairs; Grid0821 and Grid0828 capture is blocked by T-PERC-GRID0821-SPAWN-CRASH-001, which resides in the CARLA server process during sensor actor creation and is not a defect in the governed pipeline code. Due to CARLA 0.9.16 instability on the target custom map (Grid0821 crashes at map-load; see POST_FIX_STATUS.json), this measurement remains deferred. a full paired Ingolstadt comparison is classified as future work.

\paragraph{RQ4/RQ5: GNN latent gap.}
An earlier positive-only cosine-loss interpretation that reported cosine similarity 0.9989 is not used as thesis evidence because that configuration was later identified as collapsed. The governed contrastive artifacts instead indicate a non-collapsed configuration, with cross-cosine $=-0.160$, training cosine similarity $=-0.206$ (cross-cosine from \texttt{collapse\_check.json}; training cosine from \texttt{gnn\_training\_report.json}), final loss $=2.477$, and an overall verdict of \texttt{NOT\_COLLAPSED}. This corrected contrastive setup is observational only and does not support a transfer-feasibility claim, because no downstream perception transfer experiment was executed. The training objective follows the SimCLR / NT-Xent family of contrastive methods \cite{chen2020simclr}.
